\begin{entry}{Q010}{The Tsirelson bound in one identity}

\entrymeta{
  topic   = {Bell nonlocality; operator norms; commutators.},
  est     = $\sim$12 minutes,
  origin  = {Adapted. The squaring identity is the standard argument of Tsirelson (1980),
    later streamlined by Landau; the staging as a self-contained operator exercise is mine.},
  source  = {B. S. Tsirelson, \emph{Quantum generalizations of Bell's inequality},
    Lett. Math. Phys. \textbf{4}, 93 (1980).},
  prereq  = {$\opnorm{\cdot}$ denotes the operator (spectral) norm.},
  status  = open,
}

\begin{problem}
Let $A_0,A_1,B_0,B_1$ be Hermitian operators on a Hilbert space $\mathcal H$ satisfying
\[
  A_i^{2}=B_j^{2}=\id
  \qquad\text{and}\qquad
  [A_i,B_j]=0 \quad\text{for all } i,j\in\{0,1\}.
\]
(These are the \emph{dichotomic observables} of two space-like separated parties: squaring
to the identity encodes outcomes $\pm 1$, and commutation encodes locality.)

Define the \emph{CHSH operator}
\[
  \mathcal{B} \;=\; A_0B_0+A_0B_1+A_1B_0-A_1B_1 .
\]

Prove that
\[
  \opnorm{\mathcal{B}} \;\le\; 2\sqrt{2}.
\]
Conclude that $\lvert\avg{\psi}{\mathcal{B}}{\psi}\rvert\le 2\sqrt{2}$ for every state
$\ket{\psi}$, and state which step of your argument would have to fail for the value $4$
--- the algebraic maximum of a sum of four terms each bounded by $1$ --- to be reachable.
\end{problem}

% ============================================================== solutions ===

\begin{solution}
Define $$B_{\pm} = \frac{1}{\sqrt{2}}(B_0 \pm B_1),$$ thus $$B_{\pm}^2=1\pm\frac{1}{2} \{B_0,B_1\} $$ where  $\{B_0,B_1\}$ is the anti-commutator. The sum of the square of the operators is 
\begin{equation}
    B_{+}^2+B_{-}^2=2
    \label{square sum}
\end{equation}
Therefore, using \eqref{square sum} together with $A_i$ being dichotomic observables
\begin{equation}
    (A_0 + B_+)^2 +(A_1+B_-)^2= 4 + 2A_0B_+ + 2A_1B_- = 4 + \sqrt{2}\beta.
\end{equation}
The left side of the equation is clearly a positive operator, hence $\avg{\psi}{\beta}{\psi} \leq 2\sqrt{2}$ for any state $\ket{\psi}$. 

We would need all 4 coefficients to be positive ($+1$) in order to $4$ be reachable, the $-$ sign forces the bell functional to be written in terms of the variables $B_{\pm}$, that counter balances each other.


\end{solution}

\begin{takeaway}
If they are classical variables, one of them is certainly $0$. For the quantum case, you can make the operators have the same expected value (which is better than maximizing one, since it forces the other to be null), leading to a quantum bound greater than the classical one. If all coefficients are $+1$, the functional is trivial, a classical deterministic strategy where all variables acquire the value $+1$ achieves the bound, and even for the quantum case 
\begin{equation}
    ||\beta ||_{\infty} \leq ||A_0B_0||_{\infty}  + ||A_0B_1||_{\infty}  + ||A_1B_0||_{\infty}  - ||A_1B_1||_{\infty} = 4 
\end{equation}
\end{takeaway}

\attempt{
  date        = 2026-09-13,
  time        = 12min,
  difficulty  = 1,
  howitwent = It was easy since I am really familiar with the result due to my research but the problems have a lot of subtitles
}

\end{entry}